\documentclass[a4paper,reqno,10pt]{amsart}
\usepackage{a4wide,amssymb,amsmath,enumitem,tikz-cd}
\usepackage[colorlinks=true,linkcolor=blue,citecolor=blue,urlcolor=blue,pdfusetitle]{hyperref}
\setlist[enumerate,1]{label={\upshape(\arabic*)}}
\setlist[enumerate,2]{label={\upshape(\alph*)}}
\numberwithin{equation}{section}
\newtheorem{theorem}{Theorem}[section]
\newtheorem{proposition}[theorem]{Proposition}
\newtheorem{lemma}[theorem]{Lemma}
\newtheorem{corollary}[theorem]{Corollary}
\newtheorem{introthm}{Theorem}
\renewcommand{\theintrothm}{\Alph{introthm}}
\theoremstyle{definition}
\newtheorem{construction}[theorem]{Construction}
\newtheorem{example}[theorem]{Example}
\newtheorem*{organization}{Organization}
\newtheorem*{conventions}{Conventions and notation}
\newtheorem*{aiuse}{Use of AI}
\newcommand{\B}{\mathcal B}
\newcommand{\C}{\mathcal C}
\newcommand{\E}{\mathbb E}
\newcommand{\T}{\mathcal T}
\newcommand{\Q}{\mathcal Q}
\newcommand{\Rcat}{\mathcal R}
\newcommand{\Kb}{\mathsf K^{\mathrm b}}
\newcommand{\Db}{\mathsf D^{\mathrm b}}
\DeclareMathOperator{\Hom}{Hom}
\DeclareMathOperator{\Ext}{Ext}
\DeclareMathOperator{\im}{im}
\DeclareMathOperator{\coker}{coker}

\title[An extriangulated category with no triangulated embedding]
{An extriangulated category with no triangulated embedding}
\author[H.\ Enomoto]{Haruhisa Enomoto}
\address{Parakeet Inc., Japan}
\email{haruhisa.enomoto.math@gmail.com}
\subjclass[2020]{18G80, 18G25, 18E05, 03E05}
\keywords{extriangulated category, triangulated category, algebraic triangulated category, inverse limit, Aronszajn tree}
\date{\today}

\begin{document}
\begin{abstract}
Over any field, we construct an essentially small hereditary extriangulated
category with enough injectives that admits no triangulated embedding, that is,
no fully faithful extriangulated functor into a triangulated category. In
particular, not every extriangulated category is equivalent as an
extriangulated category to a full extension-closed subcategory of a
triangulated category. Nevertheless, every countable family of objects lies in
a full extension-closed subcategory equivalent as an extriangulated category to
a full extension-closed subcategory of an algebraic triangulated category. To
construct this example, we build an extriangulated category from any inverse
system of vector spaces with surjective transition maps over a nonempty
ordinal. We prove that it admits a triangulated embedding if and only if every
projection from the inverse limit is surjective. We obtain the counterexample
from such an inverse system over the countable ordinals with nonzero initial
term and vanishing inverse limit.
\end{abstract}
\maketitle

\section{Introduction}

Extriangulated categories were introduced by Nakaoka and Palu
\cite[Definition~2.12]{NP} to provide a common framework for triangulated
categories and exact categories with set-valued first extension groups
\cite[Example~2.13 and Proposition~3.22(1)]{NP}. A primary source of further
examples is any full extension-closed subcategory $\C$ of a triangulated
category $\T$. For $A,C\in\C$, each morphism $\delta:C\to A[1]$ completes to a
distinguished triangle $A\to B\to C\xrightarrow{\delta}A[1]$ with $B$
isomorphic to an object of $\C$, and assigning the equivalence class of
$A\to B\to C$ to $\delta$ equips $\C$ with an extriangulated structure
\cite[Remark~2.18]{NP}. Gorsky, Nakaoka and Palu \cite[\S1, opening
paragraph]{GNP} asked whether every extriangulated category arises in this way,
that is, is equivalent as an extriangulated category to a full extension-closed
subcategory of a triangulated category. Notice that an ambient triangulated
category supplies extra structure beyond the extriangulated axioms: applying
$\Hom$ to its distinguished triangles extends the $\Hom$--$\E$ exact sequences
to the left via negative shifts.

In this paper, we answer the question of Gorsky, Nakaoka and Palu negatively
over every field. In fact, our counterexample admits no fully faithful
extriangulated functor into any triangulated category, even though it is
hereditary and Krull--Schmidt with enough injectives, and every countable
family of its objects lies in a full extension-closed subcategory that embeds
into an algebraic triangulated category.

To state the results precisely, recall that an extriangulated category
$\mathcal A$ is equipped with a biadditive bifunctor $\E(C,A)$ of extensions,
contravariant in $C$ and covariant in $A$, together with an assignment sending
each extension class $\varepsilon\in\E(C,A)$ to a pair $A\to B\to C$ (called a
\emph{conflation}) up to isomorphism of middle terms fixing $A$ and $C$. An additive functor
$F:\mathcal A\to\T$ into a triangulated category is \emph{extriangulated} in
the sense of Bennett-Tennenhaus and Shah \cite[Definition~2.32 with $n=1$]{BS}
if there exists a natural family of group homomorphisms
\[
  \eta_{C,A}:\E(C,A)\longrightarrow\T(FC,FA[1])
\]
such that $FA\to FB\to FC\xrightarrow{\eta_{C,A}(\varepsilon)}FA[1]$ is a
distinguished triangle in $\T$ for every conflation $A\to B\to C$ realizing
$\varepsilon$. We call such a pair $(F,\eta)$ a \emph{triangulated embedding}
if $F$ is fully faithful, and a \emph{triangulated realization} if, in
addition, every $\eta_{C,A}$ is an isomorphism. If $(F,\eta)$ is a triangulated
realization, then every morphism $FC\to FA[1]$ is of the form
$\eta_{C,A}(\varepsilon)$ for some $\varepsilon\in\E(C,A)$, so the middle term
of its distinguished triangle is isomorphic to $FB$ for a conflation
$A\to B\to C$. Hence the essential image of $F$ is extension-closed in $\T$,
and $F$ is an extriangulated equivalence onto it with the induced structure.
Conversely, the inclusion of a full extension-closed subcategory with its
induced structure is a triangulated realization. A triangulated category is
\emph{algebraic} if it is triangle equivalent to the stable category of a
Frobenius exact category, and a triangulated realization into an algebraic
triangulated category is called an \emph{algebraic realization}.

For an essentially small extriangulated category $\mathcal A$, we use the
positive Yoneda extensions $\E^n$ defined by Gorsky, Nakaoka and Palu
\cite[Definition~3.1 and Remarks~3.2--3.3]{GNP} on a small skeleton. The
bifunctor $\E^n$ is the $n$-th tensor power of $\E$ over $\mathcal A$, so its
elements are sums of products of $n$ composable extension classes. By
\cite[Proposition~2.1(1)(iii) and Definition~2.3]{GNP2}, $\mathcal A$ is
\emph{hereditary} precisely when $\E^2=0$. An object $U\in\mathcal A$ is
\emph{injective} if every morphism $A\to U$ factors through $A\to B$ for every
conflation $A\to B\to C$, or equivalently if $\E(-,U)=0$ \cite[dual of
Proposition~3.24]{NP}. The category $\mathcal A$ has \emph{enough injectives}
if every object occurs as the first term of a conflation whose middle term is
injective.

Since every triangulated realization is a triangulated embedding, part~(1) of
our first main theorem answers the question of Gorsky, Nakaoka and Palu
negatively.

\begin{introthm}[= Theorem~\ref{thm:example}]\label{intro:example}
Over every field $k$ there exists an essentially small $k$-linear
extriangulated category $\C$ with the following properties.
\begin{enumerate}
\item The category $\C$ has no triangulated embedding.
\item Every countable family of objects of $\C$ is contained in a full
extension-closed subcategory having an algebraic realization.
\item The category $\C$ is hereditary and Krull--Schmidt, and has enough
injectives. It has $\aleph_1$ isomorphism classes of indecomposables.
Each indecomposable has endomorphism ring $k$, and
$\dim_k\E(C,A)\leq 1$ for all indecomposables $A,C$. Every morphism space
has at most countable dimension over $k$.
\end{enumerate}
\end{introthm}

Theorem~\ref{intro:example} also has consequences for enhancements by exact
$\infty$-categories and exact dg categories. Klemenc \cite[Theorem~1.2]{Kl}
showed that every small exact $\infty$-category embeds into its stable hull as
an extension-closed subcategory. On homotopy categories, this produces a triangulated realization of
the extriangulated structure constructed by Nakaoka and Palu
\cite[Theorem~4.22]{NP2} (see \cite[\S5.5, before Question~5.40]{GNP}); in
particular, this covers small ordinary exact categories via their nerves
\cite[\S4.3]{Kl}. Furthermore, Chen \cite[Proposition-Definition~3.15]{Ch}
characterized the essentially small extriangulated categories admitting an
algebraic realization as those equivalent, as extriangulated categories, to the
homotopy categories of exact dg categories. Consequently,
Theorem~\ref{intro:example}(1) shows that $\C$ is not equivalent as an
extriangulated category to the homotopy category of any small exact
$\infty$-category or exact dg category.

We obtain $\C$ as a member of a general family of extriangulated categories
associated to inverse systems of vector spaces. Let $I$ be a nonempty ordinal
and let $D=(D_i,\pi_{ji}:D_j\to D_i)_{i\le j\text{ in }I}$ be an inverse system
of $k$-vector spaces with surjective transition maps. Adjoining a new least
element $*$ below $0\in I$, we consider covariant representations of the
resulting poset. Write $A_i$ and $P_i$ for the interval representations
supported on $[0,i]$ and $[*,i]$, respectively, with value $k$ and identity
maps on their supports. The inclusions $A_i\to P_i$ share the common quotient
$Y$, supported at $*$, and for $i\le j$ the quotient map $P_j\to P_i$ restricts
to $A_j\to A_i$ and induces the identity on $Y$.

In Construction~\ref{con:category}, we adjoin a single object $X$ to the full
additive subcategory generated by $Y$ and the $A_i,P_i$ so that the morphism
spaces from $X$ to the $A_i$ form the inverse system $D$, while all morphisms
from $X$ to $P_i$ or $Y$ (and from these interval objects to $X$) vanish. In
particular, every composite $X\to A_i\to P_i$ is zero. We equip the resulting
additive category $\C(D)$ with the extension bifunctor pulled back from the
interval representations, so that each interval sequence $A_i\to P_i\to Y$ is a
conflation in $\C(D)$. Our second main theorem characterizes when $\C(D)$
embeds into a triangulated category.

\begin{introthm}[= Theorem~\ref{thm:criterion}]\label{intro:criterion}
For every inverse system $D=(D_i)_{i\in I}$ of vector spaces over a field,
with surjective transition maps and indexed by a nonempty ordinal $I$, the
following conditions on $\C(D)$ and $D$ are equivalent.
\begin{enumerate}
\item The category $\C(D)$ admits a triangulated embedding.
\item The projection $\varprojlim D\to D_i$ is surjective for every $i\in I$.
\item The category $\C(D)$ has an algebraic realization.
\end{enumerate}
\end{introthm}

Let us explain the idea behind Theorem~\ref{intro:criterion}. Given a
triangulated embedding $(F,\eta):\C(D)\to\T$, the group $\T(FX,FY[-1])$ serves
as the ambient group of negative first extensions from $X$ to $Y$
\cite[\S5.5]{GNP}. Applying $\T(FX,-)$ to the distinguished triangle
$FA_i\to FP_i\to FY\to FA_i[1]$ yields a surjection from $\T(FX,FY[-1])$ onto
$\T(FX,FA_i)\simeq D_i$, because every morphism $X\to A_i$ composes to zero
with $A_i\to P_i$. Since the transition maps between the interval conflations
induce the identity on $Y$, these surjections are compatible with the
transition maps $\pi_{ji}$, and therefore induce a homomorphism
$\T(FX,FY[-1])\to\varprojlim D$ whose composite with every projection
$\varprojlim D\to D_i$ is surjective. Conversely, given any vector space $W$
admitting compatible surjections onto all $D_i$, we construct in
Proposition~\ref{prop:complex} an algebraic realization of $\C(D)$ inside a
bounded homotopy category $\Kb(\Q)$ in which $W$ realizes this degree $-1$
morphism space.

To prove Theorem~\ref{intro:example}, we apply Theorem~\ref{intro:criterion} to
the inverse system of free vector spaces on the levels of an Aronszajn tree, as
studied by Dimitri\'c \cite[FC2, pp.~20--21]{Di}. By Bergman \cite[p.~6, argument
preceding Corollary~8]{Be2} (see also Lemma~\ref{lem:linear}), this inverse system
has trivial limit, whereas the tree in Construction~\ref{con:tree} has a single
root and thus $D_0=k$. Consequently, the projection $\varprojlim D\to D_0$ is
not surjective, so $\C(D)$ admits no triangulated embedding by
Theorem~\ref{intro:criterion}. On the other hand, every countable set of
countable ordinals has a countable upper bound $m$, and the transition maps
from $D_m$ surject onto all terms at indices $\le m$, yielding an algebraic
realization of the corresponding full extension-closed subcategory.

Theorem~\ref{intro:example}(2) shows that the obstruction to a triangulated
embedding involves uncountably many indecomposables. In
Proposition~\ref{prop:bounds}, we show that within the family $\C(D)$ both the
size of the index set and the dimension of the morphism spaces in
Theorem~\ref{intro:example} are minimal. Specifically, any counterexample of
the form $\C(D)$ must be indexed by an ordinal of uncountable cofinality (so
$|I|\ge\aleph_1$) and must have an infinite-dimensional morphism space
$\C(D)(X,A_i)=D_i$. We do not know whether every essentially small $k$-linear
extriangulated category with finite-dimensional morphism spaces has a
triangulated realization.

\begin{organization}
In Section~\ref{sec:intervals}, we construct the categories $\C(D)$ and their
subcategories $\C_J$ from interval representations, and we prove in
Proposition~\ref{prop:complex} that $\C_J$ admits an algebraic realization
whenever the spaces $D_i$ ($i\in J$) are compatible quotients of a single
vector space. Applying this to finite subsets verifies the extriangulated
axioms for $\C(D)$. In Section~\ref{sec:criterion}, we prove
Theorem~\ref{intro:criterion}. In Section~\ref{sec:example}, we first obtain a
quick counterexample from the inverse system of injections of countable
ordinals into $\mathbb N$ given by Bergman \cite[p.~7]{Be2}
and later by Wofsey \cite{Wo}, and then pass to an
Aronszajn tree with countable levels to prove Theorem~\ref{intro:example}.
Finally, Proposition~\ref{prop:bounds} establishes the minimality of the
cardinality and dimension bounds within the family $\C(D)$.
\end{organization}

\begin{conventions}
We fix a field $k$ and equip every ordinal with its usual well-order. For a
category $\mathcal A$, we write $\mathcal A(C,A)$ for the set of morphisms from
$C$ to $A$. Subcategories specified by collections of objects and their finite
direct sums are understood to be closed under isomorphism. We use cochain
complexes with shift convention $M[1]^n=M^{n+1}$, and write $\Kb(\mathcal A)$
for the bounded homotopy category of an additive category $\mathcal A$.
\end{conventions}

\begin{aiuse}
The author posed the problem of settling the triangulated realization question
for extriangulated categories and directed the search for a simple, conceptual
counterexample. The counterexample, embedding criterion, and proofs presented
here were developed by GPT-6.2 Astra under the author's direction. The initial
manuscript, citation check, and mathematical audits were carried out by
GPT-6.2 Astra; independent exposition reviews were performed by Claude
Opus~5.5; and the structure, exposition, and proof presentation were polished
with Gemini~4 Neon. The author directed the research and revisions and is
responsible for the mathematical content and the final manuscript.
\end{aiuse}

\section{Construction of the extriangulated category from intervals}\label{sec:intervals}

Throughout this section, let $I$ be a nonempty ordinal and let
$D=(D_i,\pi_{ji}:D_j\to D_i)_{i\le j\text{ in }I}$ be an inverse system of
$k$-vector spaces with surjective transition maps. For each subset
$J\subseteq I$, we construct the full subcategory $\C_J$ of $\C(D)$ generated
under finite direct sums by $X$, $Y$, and the interval objects $A_i,P_i$ with
$i\in J$. We first study the underlying category of interval representations,
then construct $\C_J$ and realize it by bounded complexes whenever a single
vector space maps compatibly onto all $D_i$ ($i\in J$), and finally use finite
subsets $J$ to verify the extriangulated category axioms for $\C(D)$.

Since $I$ is a nonempty ordinal, it has a least element $0$. We adjoin a new
least element $*$ below $0$ and write $I_*$ for the resulting well-ordered set.
A representation $M$ of $I_*$ is a covariant functor from $I_*$ to the category
of $k$-vector spaces, assigning a $k$-vector space $M(t)$ to each $t\in I_*$
and a linear structure map $M(s\le t):M(s)\to M(t)$ to each pair $s\le t$ in
$I_*$. These representations form an abelian category $\Rcat$ in which
exactness is checked pointwise.

For an interval $U\subseteq I_*$, let $k_U\in\Rcat$ be the interval
representation with value $k$ on $U$ and $0$ elsewhere, and with identity
structure maps between points of $U$. We work with the intervals
\[
  Y=k_{\{*\}},\qquad A_i=k_{[0,i]},\qquad P_i=k_{[*,i]}
       \qquad(i\in I).
\]
Write $b_i:A_i\to P_i$ for the canonical inclusion and $c_i:P_i\to Y$ for the
canonical quotient. For $i\leq j$ in $I$, there are canonical quotient maps
$a_{ji}:A_j\to A_i$ and $p_{ji}:P_j\to P_i$, which fit into a morphism of short
exact sequences
\begin{equation}\label{eq:rows}
\begin{tikzcd}[column sep=1.8em]
0\ar[r]&A_j\ar[r,"b_j"]\ar[d,"a_{ji}"']&P_j\ar[r,"c_j"]\ar[d,"p_{ji}"]&
Y\ar[r]\ar[d,equal]&0\\
0\ar[r]&A_i\ar[r,"b_i"']&P_i\ar[r,"c_i"']&Y\ar[r]&0 .
\end{tikzcd}
\end{equation}
For $J\subseteq I$, let $\B_J$ be the full subcategory of $\Rcat$ consisting of
finite direct sums of $Y$ and the $A_i,P_i$ with $i\in J$. We begin by showing
that each $\B_J$ is extension-closed in $\Rcat$ and has enough injectives.

\begin{lemma}\label{lem:intervals}
For every $J\subseteq I$, the subcategory $\B_J$ is extension-closed in
$\Rcat$. With its induced exact structure, it has enough injectives.
The objects $P_i$ and $Y$ are injective in $\Rcat$, and
$0\to A_i\xrightarrow{b_i}P_i\xrightarrow{c_i}Y\to0$ is an injective
resolution for each $i\in J$.
\end{lemma}
\begin{proof}
Botnan and Crawley-Boevey \cite[Lemma~2.1 and Theorem~1.2]{BCB} prove
injectivity for interval representations supported on directed ideals and
interval decomposition for pointwise finite-dimensional representations of an
arbitrary total order. For the principal ideals needed here, we give a direct
two-step argument: evaluation yields injectivity, and splitting on $I$
determines the middle term of any extension in $\B_J$.

\medskip\noindent\textbf{Step 1: Injective intervals.}
For any $t\in I_*$ and $M\in\Rcat$, evaluation at $t$ yields a natural
isomorphism
\[
  \Rcat(M,k_{[*,t]})\ \simeq\ \Hom_k(M(t),k).
\]
Indeed, a morphism $\varphi:M\to k_{[*,t]}$ vanishes at every point strictly
above $t$, and for $s\le t$ its component $\varphi_s:M(s)\to k$ is uniquely
determined by $\varphi_t:M(t)\to k$ via $\varphi_s=\varphi_t\circ M(s\le t)$;
conversely, every linear functional on $M(t)$ defines a morphism in this way.
Since evaluation at $t$ and vector space duality are exact, $k_{[*,t]}$ is
injective in $\Rcat$. Taking $t=i$ and $t=*$ shows that $P_i$ and $Y$ are
injective in $\Rcat$. The same argument applied to $I$ in place of $I_*$ shows
that $k_{[0,i]}$ is injective in the category of representations of $I$.

\medskip\noindent\textbf{Step 2: Closure of $\B_J$ under extensions.}
Consider a short exact sequence in $\Rcat$ with outer terms in $\B_J$ and
middle term $M$. On $I$, we have $Y|_I=0$ and $A_i|_I=P_i|_I=k_{[0,i]}$. Hence
the restriction of the first term to $I$ is a finite direct sum of the
injective representations $k_{[0,i]}$, so the restricted short exact sequence
splits. It follows that $M|_I$ is a finite direct sum of intervals $k_{[0,i]}$
whose endpoints belong to the finite subset of $J$ occurring in the outer
terms. List the distinct endpoints occurring in $M|_I$ in increasing order as
$i_1<\dots<i_r$.

If $r=0$, then $M|_I=0$, so $M$ is supported at $*$ and hence
$M\simeq Y^{\dim_k M(*)}\in\B_J$ (as $M(*)$ is finite-dimensional). Assume now
that $r\ge 1$, and write $V=M(0)$ and $d=M(*\le 0):M(*)\to V$. Define a flag
$0=V_0\subset V_1\subset\dots\subset V_r=V$ by setting
\[
  V_0=0,\qquad V_s=\ker\bigl(M(0\le i_{s+1}):V\to M(i_{s+1})\bigr)\ \ (0<s<r),\qquad V_r=V.
\]
Under the decomposition of $M|_I$ into intervals $k_{[0,i_u]}$, the subspace
$V_s$ is the direct sum of the values at $0$ of the summands with $u\le s$,
since $k_{[0,i_u]}(t)$ equals $k$ for $t\le i_u$ and $0$ for $t>i_u$.
Consequently, for every $t\in I$, if $s$ is the number of indices
$u\in\{1,\dots,r\}$ with $i_u<t$, then the structure map $M(0\le t):V\to M(t)$
is surjective with kernel $V_s$, inducing an isomorphism $V/V_s\simeq M(t)$ under
which the structure maps on $I$ correspond to the canonical quotient maps
$V/V_s\to V/V_{s'}$.

Set $W_s=\im d\cap V_s$ for $0\le s\le r$, and put
$m_s=\dim_k(W_s/W_{s-1})$ and $n_s=\dim_k(V_s/V_{s-1})-m_s$ for $1\le s\le r$.
For each $s=1,\dots,r$, we choose vectors $v_{s,1},\dots,v_{s,m_s}\in W_s$
whose cosets form a basis of $W_s/W_{s-1}$. Because
$W_s\cap V_{s-1}=W_{s-1}$, their cosets in $V_s/V_{s-1}$ are linearly
independent, so we may extend them to a basis of $V_s/V_{s-1}$ by adjoining
the cosets of vectors $w_{s,1},\dots,w_{s,n_s}\in V_s$. Together, all
$v_{s,\alpha}$ and $w_{s,\beta}$ form a basis of $V$ adapted to the flag
$V_0\subset\dots\subset V_r$. For each $v_{s,\alpha}\in\im d$, choose a
preimage $\tilde v_{s,\alpha}\in M(*)$ with $d(\tilde v_{s,\alpha})=v_{s,\alpha}$,
and complete these lifts to a basis of $M(*)$ by adjoining a basis of $\ker d$.
Each basis vector of $\ker d$ generates a subrepresentation isomorphic to $Y$;
each $\tilde v_{s,\alpha}$ generates a subrepresentation isomorphic to
$P_{i_s}$ (since $v_{s,\alpha}\in V_s\setminus V_{s-1}$); and each
$w_{s,\beta}$ (with zero component at $*$) generates a subrepresentation
isomorphic to $A_{i_s}$. Since these subrepresentations span $M(t)$ linearly
independently at every $t\in I_*$ and $M(*)$ is finite-dimensional, we obtain
\[
 M\simeq Y^{\dim_k\ker d}\oplus
 \bigoplus_{s=1}^r\bigl(P_{i_s}^{m_s}\oplus A_{i_s}^{n_s}\bigr).
\]
Thus $M\in\B_J$, proving that $\B_J$ is extension-closed in $\Rcat$. Finally,
since $P_i$ and $Y$ are injective in $\Rcat$ (hence in $\B_J$), the sequences
$0\to A_i\xrightarrow{b_i}P_i\xrightarrow{c_i}Y\to0$ are injective resolutions in
$\B_J$; taking finite direct sums with the identity maps on $P_i$ and $Y$ shows
that $\B_J$ has enough injectives.
\end{proof}

Before adjoining $X$ to $\B_J$, we record a basis for each morphism space
between the indecomposable objects $A_j,P_j,Y$ of $\B_J$. In the table below,
each nonzero entry spans the corresponding morphism space
$\B_J(\text{source},\text{target})$, and an entry marked with the condition
$i\le j$ is understood to be $0$ when $i>j$:
\begin{equation}\label{eq:hom}
\begin{array}{c|ccc}
 \text{source}\ \backslash\ \text{target}& A_i&P_i&Y\\ \hline
 A_j & a_{ji}\ (i\le j)& b_i a_{ji}\ (i\le j)&0\\
 P_j & 0 & p_{ji}\ (i\le j) & c_j\\
 Y   & 0 & 0 & 1_Y
\end{array}
\end{equation}
Indeed, by Step~1 of the proof of Lemma~\ref{lem:intervals}, we have
$\Rcat(M,P_i)\simeq\Hom_k(M(i),k)$ and $\Rcat(M,Y)\simeq\Hom_k(M(*),k)$ via
evaluation at $i$ and $*$, respectively. In particular, $\Rcat(A_j,P_i)$ is
one-dimensional with basis $b_i a_{ji}$ when $i\le j$ (since $b_i a_{ji}$ acts
as the identity at $i$) and vanishes when $i>j$; the remaining entries in the
columns for $P_i$ and $Y$ follow in the same way. Finally, because
$0\to A_i\xrightarrow{b_i}P_i\xrightarrow{c_i}Y\to0$ is exact, morphisms
$M\to A_i$ correspond to morphisms $M\to P_i$ whose composite with $c_i$
vanishes, yielding the first column of \eqref{eq:hom}.

\begin{construction}\label{con:category}
Let $D=(D_i,\pi_{ji}:D_j\to D_i)_{i\leq j\text{ in }I}$ be an inverse system of
$k$-vector spaces with surjective transition maps. For $J\subseteq I$, we form
a $k$-linear additive category $\C_J$ by adjoining an object $X$ to $\B_J$ and
taking finite direct sums, with morphism spaces
\begin{equation}\label{eq:X}
  \C_J(X,X)=k,\qquad \C_J(B,X)=0,\qquad
  \C_J(X,A_i)=D_i,\qquad \C_J(X,P_i)=0,\qquad \C_J(X,Y)=0
\end{equation}
for all $B\in\B_J$ and $i\in J$. For $i\le j$ in $J$, we define postcomposition
with $a_{ji}:A_j\to A_i$ on $\C_J(X,A_j)=D_j$ to be $\pi_{ji}:D_j\to D_i$.
Together with \eqref{eq:hom}, $k$-bilinearity, and the action of
$\C_J(X,X)=k$ by scalar multiplication, this determines composition on all
objects of $\C_J$. To check associativity on morphisms out of $X$, we only need
to verify composites $X\to A_j\to A_i\to A_h$ with $h\le i\le j$, since every
morphism from $X$ to $P_i$ or $Y$ is zero; on the generators $a_{ji}$ and
$a_{ih}$, associativity is precisely the identity $\pi_{ih}\pi_{ji}=\pi_{jh}$.

Thus the underlying additive category of $\C_J$ is the additive closure of the
one-point extension of $\B_J$ by the covariant module $B\mapsto\C_J(X,B)$, in
the sense of Herscovich and Solotar \cite[\S3]{HeS}. We write $\C(D)=\C_I$.
Every object of $\C_J$ has the form $B\oplus X^a$ with $B\in\B_J$ and
$a\ge 0$. Sending $B\oplus X^a$ to $B$ and each morphism to its component
between the $\B_J$-summands defines a $k$-linear functor $\rho:\C_J\to\B_J$
with $\rho X=0$; it respects composition because $\C_J(B,X)=0$ for all
$B\in\B_J$. We define a biadditive bifunctor $\E_J$ on $\C_J$ by
\[
    \E_J(C,A)=\Ext^1_{\B_J}(\rho C,\rho A)
\]
on objects and morphisms, and write $\E=\E_I$ for $\C(D)$. In particular, every
morphism $u:X\to A_i$ acts as zero on $\E_J$ because $\rho u=0$. For objects
$A=\rho A\oplus X^a$ and $C=\rho C\oplus X^c$ of $\C_J$, let a class
$\xi\in\E_J(C,A)$ be represented by a short exact sequence
$0\to\rho A\to M\to\rho C\to0$ in $\B_J$. We assign to $\xi$ the equivalence
class (up to isomorphism of middle terms fixing $A$ and $C$) of the direct sum
of $\rho A\to M\to\rho C$ and the split sequence
$X^a\to X^a\oplus X^c\to X^c$.
\end{construction}

To realize the morphisms $X\to A_i$ as chain maps modulo homotopy, we realize
the spaces $D_i$ ($i\in J$) as compatible quotients of a single vector space.

\begin{proposition}\label{prop:complex}
Suppose that a vector space $W$ admits surjections $q_i:W\to D_i$ for
$i\in J$ satisfying $\pi_{ji}q_j=q_i$ for all $i\le j$ in $J$.
Then the data of Construction~\ref{con:category} equip $\C_J$ with an
extriangulated structure admitting an algebraic realization
$(F,\eta):\C_J\to\Kb(\Q)$ for some additive category $\Q$.
\end{proposition}
\begin{proof}
We divide the proof into three steps: first we realize $\B_J$ by bounded
complexes of injectives, next we enlarge the category of terms to realize $X$
and construct a fully faithful functor $F:\C_J\to\Kb(\Q)$, and finally we
identify the extension groups and verify the algebraic realization.

\medskip\noindent\textbf{Step 1: Realizing $\B_J$ in $\Kb(\Q_0)$.}
Let $\Q_0$ be the full subcategory of $\Rcat$ consisting of finite direct sums
of $Y$ and the $P_i$ ($i\in J$). By Lemma~\ref{lem:intervals}, every object of
$\Q_0$ is injective in $\Rcat$. We represent $Y$ and each $P_i$ as stalk
complexes concentrated in degree $0$, and each $A_i$ by the two-term complex
\[
    H_i=(P_i\xrightarrow{c_i}Y)\qquad\text{in degrees }0,1.
\]
The inclusion of the stalk $A_i$ into $H_i$ is a quasi-isomorphism in
$\Db(\Rcat)$. Extending additively over finite direct sums assigns to each
$B\in\B_J$ a complex $H_B\in\Kb(\Q_0)$ together with a quasi-isomorphism
$B\to H_B$. By \cite[Tag 05TG]{Stacks}, morphisms in $\Db(\Rcat)$ into a
bounded complex of injectives are computed modulo chain homotopy, so the
canonical functor $\Kb(\Q_0)\to\Db(\Rcat)$ is fully faithful. Because
$\B_J\hookrightarrow\Rcat\hookrightarrow\Db(\Rcat)$ is also fully faithful,
$B\mapsto H_B$ defines a fully faithful $k$-linear functor
$\B_J\to\Kb(\Q_0)$. Moreover, for $B,C\in\B_J$ we obtain natural isomorphisms
\[
  \Kb(\Q_0)(H_C,H_B[1])
     \ \simeq\ \Db(\Rcat)(C,B[1])
     \ \simeq\ \Ext^1_{\Rcat}(C,B)
     \ \simeq\ \Ext^1_{\B_J}(C,B),
\]
where the last isomorphism holds because $\B_J$ is extension-closed in $\Rcat$
by Lemma~\ref{lem:intervals}. Given a short exact sequence
$0\to B\to M\to C\to0$ in $\B_J$, let $H_C\to H_B[1]$ be the morphism in
$\Kb(\Q_0)$ corresponding to its class under this isomorphism, and complete it
to a distinguished triangle $H_B\to E\to H_C\to H_B[1]$ in $\Kb(\Q_0)$. In
$\Db(\Rcat)$, this triangle is isomorphic to the distinguished triangle
$B\to M\to C\to B[1]$ by an isomorphism restricting to the chosen
quasi-isomorphisms on $B$ and $C$. Since $M\in\B_J$, full faithfulness of
$\Kb(\Q_0)\to\Db(\Rcat)$ yields an isomorphism $E\simeq H_M$ in $\Kb(\Q_0)$
identifying the first two arrows of the triangle with the images of
$B\to M\to C$.

\medskip\noindent\textbf{Step 2: Adjoining $L$ and defining $F$.}
Set $K_i=\ker q_i$ for each $i\in J$. To represent morphisms $X\to A_i$ by
homotopy classes in $W/K_i\simeq D_i$, we arrange that $W$ is the space of
chain maps and $K_i$ its subspace of null-homotopic maps. As in
Construction~\ref{con:category}, let $\Q$ be the additive closure of the
one-point extension obtained by adjoining a new object $L$ to $\Q_0$, with
\[
  \Q(L,L)=k,\qquad \Q(Q,L)=0,\qquad \Q(L,P_i)=K_i,\qquad \Q(L,Y)=W
\]
for all $Q\in\Q_0$ and $i\in J$. For $i\le j$ in $J$, the identity
$q_i=\pi_{ji}q_j$ implies $K_j\subseteq K_i$. We define postcomposition of
morphisms out of $L$ with $p_{ji}:P_j\to P_i$ and $c_i:P_i\to Y$ to be the
inclusions $K_j\hookrightarrow K_i$ and $K_i\hookrightarrow W$, respectively.
By \eqref{eq:hom}, the composition relations among generators of $\Q_0$ are
$c_i p_{ji}=c_j$ and $p_{ih}p_{ji}=p_{jh}$ ($h\le i\le j$ in $J$), which act on
$K_j$ by the inclusions $K_j\hookrightarrow W$ and $K_j\hookrightarrow K_h$, so
composition is associative. Since chain maps and homotopies between complexes
with terms in $\Q_0$ are unchanged in $\Q$, the inclusion
$\Kb(\Q_0)\hookrightarrow\Kb(\Q)$ is fully faithful.

We represent $X$ by the stalk complex $L[-1]$, concentrated in degree $1$.
Since $L[-1]$ and $Y[-1]$ are stalks in degree $1$, we have
$\Kb(\Q)(L[-1],Y[-1])=\Q(L,Y)=W$. Similarly, a chain map $L[-1]\to H_i$ is
given simply by a morphism $L\to Y$ in degree $1$, i.e., an element of $W$, and
it is null-homotopic precisely when it equals $c_i h$ for some
$h\in\Q(L,P_i)=K_i$. Hence $q_i$ induces an isomorphism
\begin{equation}\label{eq:quotient}
    \Kb(\Q)(L[-1],H_i)=W/K_i\xrightarrow[\ \sim\ ]{\bar q_i} D_i .
\end{equation}
We define $F$ on $\C_J(X,A_i)=D_i$ to be $\bar q_i^{-1}$, sending $u\in D_i$ to
the homotopy class of any $w\in W$ with $q_i(w)=u$. By \eqref{eq:rows}, the
chain map $H_j\to H_i$ representing $a_{ji}$ has components $(p_{ji},1_Y)$, so
postcomposition with it on \eqref{eq:quotient} is the canonical projection
$W/K_j\to W/K_i$, which corresponds under $\bar q_j$ and $\bar q_i$ to
$\pi_{ji}$ because $\pi_{ji}q_j=q_i$. Moreover, $\Kb(\Q)(L[-1],P_i)=0$ and
$\Kb(\Q)(L[-1],Y)=0$ for degree reasons; $\Kb(\Q)(H_B,L[-1])=0$ for every
$B\in\B_J$ because $\Q(Q,L)=0$ for all $Q\in\Q_0$; and
$\Kb(\Q)(L[-1],L[-1])=\Q(L,L)=k$. Extending $B\mapsto H_B$ ($B\in\B_J$) and
$X\mapsto L[-1]$ additively therefore defines a fully faithful $k$-linear
functor $F:\C_J\to\Kb(\Q)$.

\medskip\noindent\textbf{Step 3: Identifying extensions and verifying the realization.}
For every $B\in\B_J$, the morphism spaces
\[
 \Kb(\Q)(L[-1],H_B[1]),\qquad
 \Kb(\Q)(H_B,L),\qquad
 \Kb(\Q)(L[-1],L)
\]
all vanish: $H_B[1]$ is concentrated in degrees $-1$ and $0$ (whereas $L[-1]$
is in degree $1$), the terms of $H_B$ admit no nonzero morphisms to $L$, and
$L[-1]$ and $L$ are stalks in distinct degrees. Consequently, for any
$C,A\in\C_J$, every morphism $FC\to FA[1]$ is uniquely determined by its
component $H_{\rho C}\to H_{\rho A}[1]$. Combining this with the isomorphisms of
Step~1 gives a group isomorphism
\[
 \eta_{C,A}:\E_J(C,A)=\Ext^1_{\B_J}(\rho C,\rho A)
      \xrightarrow{\ \sim\ }\Kb(\Q)(FC,FA[1]).
\]
For morphisms $f:C'\to C$, $g:A\to A'$, and any $\xi\in\E_J(C,A)$, we have
\[
 (Fg)[1]\,\eta_{C,A}(\xi)\,(Ff)
     =\eta_{C',A'}\bigl((\rho g)_*(\rho f)^*\xi\bigr).
\]
Indeed, writing $A=\rho A\oplus X^a$, $C=\rho C\oplus X^c$, and similarly for
$A'$ and $C'$, the morphism $\eta_{C,A}(\xi)$ has only its
$(H_{\rho C},H_{\rho A}[1])$-component nonzero. Since
$\Kb(\Q)(L[-1],H_{\rho A}[1])=0$, the component of $\eta_{C,A}(\xi)(Ff)$ out of
$F(X^{c'})$ vanishes; similarly, $\eta_{C,A}(\xi)$ factors through
$H_{\rho A}[1]$, and since $\Kb(\Q)(H_{\rho A}[1],L)=0$ (as $\Q(Q,L)=0$ for all
$Q\in\Q_0$), the component of $(Fg)[1]\eta_{C,A}(\xi)$ into $F(X^{a'})[1]$
vanishes as well. On the remaining component $H_{\rho C'}\to H_{\rho A'}[1]$,
the identity follows from the naturality of the isomorphisms in Step~1.

Because each $\eta_{C,A}$ is surjective, every morphism $FC\to FA[1]$ equals
$\eta_{C,A}(\xi)$ for some $\xi\in\E_J(C,A)$. Represent $\xi$ by a short exact
sequence $0\to\rho A\to M\to\rho C\to0$ in $\B_J$. By Step~1, the image
$H_{\rho A}\to H_M\to H_{\rho C}$ completes to a distinguished triangle in
$\Kb(\Q)$ with connecting morphism corresponding to $\xi$. Taking the direct
sum with the split distinguished triangle
$F(X^a)\to F(X^a\oplus X^c)\to F(X^c)\xrightarrow{0}F(X^a)[1]$ produces a
distinguished triangle with connecting morphism $\eta_{C,A}(\xi)$ whose first
two arrows are the image under $F$ of the pair assigned to $\xi$ in
Construction~\ref{con:category}. Since every distinguished triangle with
connecting morphism $\eta_{C,A}(\xi)$ is isomorphic to this one fixing $FA$ and
$FC$, the essential image of $F$ is extension-closed in $\Kb(\Q)$, and $\eta$
identifies its induced extriangulated structure \cite[Remark~2.18]{NP} with
that of Construction~\ref{con:category}.

Finally, with the degreewise split exact structure, the category of bounded
complexes over $\Q$ is a Frobenius exact category whose projective-injective
objects are the contractible complexes. Its stable category is $\Kb(\Q)$
because the morphisms factoring through contractible complexes are precisely
the null-homotopic chain maps (see Keller \cite[\S2.3]{Ke}). Hence $(F,\eta)$
is an algebraic realization of $\C_J$.
\end{proof}

We can now verify the extriangulated category axioms for $\C(D)$ without
assuming the existence of a single space mapping compatibly onto all $D_i$
($i\in I$).

\begin{corollary}\label{cor:axioms}
Construction~\ref{con:category} yields an essentially small $k$-linear
extriangulated category $\C(D)$. Each $\C_J$ ($J\subseteq I$) is a full
extension-closed subcategory with its induced structure. If $J$ has a largest
element, then $\C_J$ admits an algebraic realization.
\end{corollary}
\begin{proof}
If $J\subseteq I$ has a largest element $m$, then applying
Proposition~\ref{prop:complex} with $W=D_m$ and $q_i=\pi_{mi}$ ($i\in J$)
provides an algebraic realization of $\C_J$; in particular, every such $\C_J$
is an extriangulated category.

Whenever $J\subseteq J'\subseteq I$, the inclusion $\C_J\hookrightarrow\C_{J'}$
is full and identifies $\E_J(C,A)$ with $\E_{J'}(C,A)$ for all $A,C\in\C_J$,
since both are naturally isomorphic to $\Ext^1_{\Rcat}(\rho C,\rho A)$ by
Lemma~\ref{lem:intervals}; moreover, it preserves the pairs assigned to
extension classes in Construction~\ref{con:category}.

The bifunctor $\E$ on $\C(D)$ is biadditive by construction. Each of the
remaining extriangulated axioms in Nakaoka and Palu \cite[Definition~2.12]{NP}
involves only finitely many objects of $\C(D)$: the realization of the zero
class and of direct sums of classes, the lifting of morphisms of extensions to
morphisms of conflations, the completion of commuting squares between
conflations, the composition diagram for two conflations, and the dual
statements. Choose a finite
nonempty subset $J\subseteq I$ containing all indices $i\in I$ appearing in the
summands of those objects. Since $J$ has a largest element, $\C_J$ already
satisfies all of these axioms. By fullness of $\C_J\hookrightarrow\C(D)$ and the
compatibility of extension groups and assigned pairs, every required diagram
and commutativity identity holds in $\C(D)$ as well.

For an arbitrary subset $J\subseteq I$, every extension class between objects
of $\C_J$ has a representative with middle term in $\C_J$, by
Lemma~\ref{lem:intervals} on the $\B_J$-summands and the split pair on the
$X$-summands; hence $\C_J$ is extension-closed in $\C(D)$ and inherits the
structure of Construction~\ref{con:category}. Finally, since $I$ is a set, the
finite direct sums of $X$, $Y$, and the $A_i,P_i$ ($i\in I$) form a set of
representatives for the isomorphism classes of objects of $\C(D)$, so $\C(D)$
is essentially small.
\end{proof}

\section{Triangulated embeddings and projections from the inverse limit}\label{sec:criterion}

Throughout this section, let $I$ be a nonempty ordinal and let
$D=(D_i,\pi_{ji})_{i\le j\text{ in }I}$ be an inverse system of $k$-vector
spaces with surjective transition maps, as in Construction~\ref{con:category}.
We now prove Theorem~\ref{intro:criterion}, showing that a triangulated
embedding of $\C(D)$ requires every projection $\varprojlim D\to D_i$ to be
surjective, and therefore (by Proposition~\ref{prop:complex}) already implies
the existence of an algebraic realization.

The essential mechanism is to lift a morphism $u:X\to A_i$ to a compatible
family in $\varprojlim D$ via a single degree $-1$ morphism in the target
triangulated category. For a single conflation in a triangulated target, the
existence of such a lift comes from exactness at the negative first connecting
map of Adachi, Enomoto and Tsukamoto \cite[Definition~2.3 (NE2) and
Example~2.4]{AET}. Intrinsically,
Gorsky, Nakaoka and Palu \cite[Definition~5.1(1), Remark~5.2 and
Definition~5.3]{GNP} describe a negative first extension in
$\E_{\mathrm I}^{-1}(X,Y)$ as a natural transformation
$\E(Y,-)\to\C(D)(X,-)$, whose connecting map for a conflation evaluates the
transformation at its extension class. Their Example~5.25 exhibits a nonzero
morphism to the first term of a conflation whose composite to the middle term
vanishes. Under a triangulated embedding $(F,\eta):\C(D)\to\T$, a
single morphism $v:FX[1]\to FY$ can be composed with $\eta_{Y,A_j}(\varepsilon_j)$
for every $j\in I$, and by naturality of $\eta$ the resulting morphisms form a
compatible family.

For each $i\in I$, let $\varepsilon_i\in\E(Y,A_i)$ be the extension class of
the interval conflation $A_i\xrightarrow{b_i}P_i\xrightarrow{c_i}Y$ in
\eqref{eq:rows}. For $i\le j$ in $I$, the morphism of short exact sequences
\eqref{eq:rows} induces an isomorphism from the pushout of the upper row along
$a_{ji}$ to the lower row fixing $A_i$ and $Y$, and hence
\begin{equation}\label{eq:push}
    (a_{ji})_*\varepsilon_j=\varepsilon_i\qquad (i\le j).
\end{equation}

\begin{theorem}\label{thm:criterion}
Let $I$ be a nonempty ordinal and let $D=(D_i)_{i\in I}$ be an inverse system
of $k$-vector spaces on $I$ whose transition maps are surjective. The following
conditions are equivalent.
\begin{enumerate}
\item The category $\C(D)$ admits a triangulated embedding.
\item For each $i\in I$, the projection $\varprojlim D\to D_i$ is surjective.
\item The category $\C(D)$ has an algebraic realization.
\end{enumerate}
\end{theorem}
\begin{proof}
$(1)\Rightarrow(2)$:
Let $(F,\eta):\C(D)\to\T$ be a triangulated embedding. For each $i\in I$, write
$\delta_i=\eta_{Y,A_i}(\varepsilon_i)$, so that the conflation
$A_i\xrightarrow{b_i}P_i\xrightarrow{c_i}Y$ maps under $F$ to a distinguished
triangle
\[
  FA_i\xrightarrow{Fb_i}FP_i\xrightarrow{Fc_i}FY\xrightarrow{\delta_i}FA_i[1]
\]
in $\T$. For $i\le j$ in $I$, naturality of $\eta$ together with
\eqref{eq:push} yields the commutative square
\begin{equation}\label{eq:triangles}
\begin{tikzcd}[column sep=2.6em]
 FY\ar[r,"\delta_j"]\ar[d,equal]&FA_j[1]\ar[d,"{(F a_{ji})[1]}"]\\
 FY\ar[r,"\delta_i"']&FA_i[1].
\end{tikzcd}
\end{equation}

Fix $i\in I$ and $u\in D_i=\C(D)(X,A_i)$. Since $\C(D)(X,P_i)=0$, we have
$b_i u=0$, and therefore $(Fb_i)[1]\circ(Fu)[1]=0$. Applying $\T(FX[1],-)$ to
the distinguished triangle for $A_i\to P_i\to Y$ gives the exact sequence
\[
 \T(FX[1],FY)\xrightarrow{(\delta_i)_*}\T(FX[1],FA_i[1])
       \xrightarrow{((F b_i)[1])_*}\T(FX[1],FP_i[1]).
\]
Hence there exists a morphism $v:FX[1]\to FY$ such that $\delta_i v=(Fu)[1]$.

For each $j\in I$, since $F$ is fully faithful and $[1]$ is an autoequivalence
of $\T$, there is a unique morphism $u_j:X\to A_j$ in $\C(D)$ satisfying
$(F u_j)[1]=\delta_j v$. In particular, $(Fu_i)[1]=\delta_i v=(Fu)[1]$, so
$u_i=u$. For $h\le j$ in $I$, the square \eqref{eq:triangles} (with $h$ in place
of $i$) yields
\[
    (F(a_{jh}u_j))[1]=(F a_{jh})[1]\,(F u_j)[1]
                     =(F a_{jh})[1]\,\delta_j v
                     =\delta_h v=(F u_h)[1].
\]
Since $F$ is faithful and $[1]$ is an autoequivalence, it follows that
$\pi_{jh}(u_j)=a_{jh}u_j=u_h$. Thus $(u_j)_{j\in I}$ is an element of
$\varprojlim D$ with $i$-th component $u$.

\medskip\noindent
$(2)\Rightarrow(3)$:
Apply Proposition~\ref{prop:complex} with $J=I$, $W=\varprojlim D$, and $q_i$
the $i$-th canonical projection $\varprojlim D\to D_i$.

\medskip\noindent
$(3)\Rightarrow(1)$:
Every algebraic realization is a triangulated embedding by definition.
\end{proof}

\section{A counterexample indexed by the countable ordinals}\label{sec:example}

By Theorem~\ref{thm:criterion}, to produce an extriangulated category $\C(D)$
with no triangulated embedding, it suffices to exhibit an inverse system $D$ of
$k$-vector spaces with surjective transition maps over a nonempty ordinal $I$
such that $\varprojlim D=0$ while $D_0\ne 0$. We construct such systems over
$\omega_1$, the first uncountable ordinal, by linearizing inverse systems of
nonempty sets with surjective transition maps and empty inverse limit. For a
set $S$, write $k^{(S)}$ for the free $k$-vector space with basis $S$. Bergman
\cite[Lemma~3]{Be} establishes the following result for direct sums of copies
of a fixed module, attributing the case of vector spaces over a field and its
finite support argument to Higman and Stone \cite[\S3]{HS}.

\begin{lemma}[{\cite[Lemma~3]{Be} with ring and module $k$}]\label{lem:linear}
Let $(S_i)_{i<\omega_1}$ be an inverse system of nonempty sets with
surjective transition maps and empty inverse limit. Then the induced
linear maps between the spaces $k^{(S_i)}$ are surjective and
$\varprojlim_{i<\omega_1} k^{(S_i)}=0$.
\end{lemma}

We first record a direct counterexample obtained from injections of countable
ordinals into $\mathbb N$.

\begin{example}\label{ex:injections}
We use the inverse system of sets in Bergman \cite[p.~7, Example~1 in
the proof of Lemma~10]{Be2}, also given in Wofsey's second example
\cite{Wo}. A related system indexed by the finite subsets of $\mathbb R$ was given earlier
by Waterhouse \cite{Wa}. For each countable ordinal $i<\omega_1$, set
\[
 S_i=\{f:i\hookrightarrow\mathbb N\mid\mathbb N\setminus\im f
                                    \text{ is infinite}\}.
\]
For $i\le j<\omega_1$, let the transition map $S_j\to S_i$ be restriction to
$i$. To see that it is surjective, let $f\in S_i$ and partition the infinite
set $\mathbb N\setminus\im f$ into two infinite subsets. Since $j\setminus i$
is countable, we may extend $f$ to an injection $j\hookrightarrow\mathbb N$
sending $j\setminus i$ into one of these subsets, leaving the other subset
outside the image. A compatible family in $\varprojlim S_i$ would have as its
union an injection $\omega_1\hookrightarrow\mathbb N$, which is impossible, so
$\varprojlim S_i=\emptyset$. Setting $D'_i=k^{(S_i)}$, Lemma~\ref{lem:linear}
gives $\varprojlim D'=0$, whereas $S_0$ consists solely of the empty function
and thus $D'_0=k$. Therefore $\C(D')$ has no triangulated embedding by
Theorem~\ref{thm:criterion}.
\end{example}

In Example~\ref{ex:injections}, the morphism space $\C(D')(X,A_\omega)=D'_\omega$
already has uncountable dimension, since $S_\omega$ contains all bijections
from $\omega=\mathbb N$ onto the even natural numbers. To obtain morphism
spaces of at most countable dimension, we replace $S_i$ by countable subsets
closed under restriction using an Aronszajn tree. Recall that a \emph{tree} is
a partially ordered set in which the strict predecessors of each node are well
ordered; its \emph{$i$-th level} consists of the nodes whose set of strict
predecessors has order type $i$. An \emph{Aronszajn tree} is a tree of height
$\omega_1$ whose $i$-th level is nonempty and countable for every $i<\omega_1$,
and which has no chain meeting every level. Aronszajn's construction is
recorded by Kurepa \cite[Th\'eor\`eme~6, p.~96, footnote~3]{Ku}.

\begin{construction}[{\cite[proof of Theorem~9.16, (9.10), pp.~116--117]{Je}}]\label{con:tree}
In the Aronszajn tree constructed by Jech \cite[proof of Theorem~9.16,
pp.~116--117]{Je}, the nodes are bounded strictly increasing sequences of
rational numbers, ordered by extension, so a node on level $i$ is a function
with domain $i$. Fixing a bijection $\mathbb Q\to\mathbb N$, since the range of
each bounded sequence omits infinitely many rationals above its upper bound, we
may identify the $i$-th level with a countable subset $N_i\subseteq S_i$ of
Example~\ref{ex:injections}. The level $N_0$ consists solely of the empty
sequence, and condition~(9.10) of \cite[pp.~116--117]{Je} ensures that every
node extends to every higher level. For $i\le j<\omega_1$, restriction to $i$
therefore defines a surjective transition map $N_j\to N_i$ (the identity when
$i=j$, and sending each node to its predecessor on level $i$ when $i<j$).
Because a compatible family would be a chain meeting every level, we have
$\varprojlim N_i=\emptyset$. Following Dimitri\'c \cite[FC2, pp.~20--21]{Di},
we set $D_i=k^{(N_i)}$, with linear transition maps induced by these
restrictions $N_j\to N_i$ (sending each basis element to its predecessor for
$i<j$). By Lemma~\ref{lem:linear}, these
linear maps are surjective and $\varprojlim D=0$. Moreover, $D_0=k$ and
$\dim_k D_i=|N_i|\le\aleph_0$ for all $i<\omega_1$.
\end{construction}

To deduce Theorem~\ref{intro:example}(3) for this tree system, we next
determine the indecomposable objects and the extension spaces of $\C(D)$ for an
arbitrary system $D$.

\begin{proposition}\label{prop:size}
For any inverse system $D=(D_i)_{i\in I}$ as in Construction~\ref{con:category},
the category $\C(D)$ is idempotent complete and Krull--Schmidt. A complete list
of pairwise nonisomorphic indecomposables is
\[
      X,\quad Y,\quad A_i,\quad P_i\qquad(i\in I).
\]
Their endomorphism rings are $k$. For $A,C$ on this list, the only nonzero
spaces $\E(C,A)$ are
\begin{equation}\label{eq:extensions}
   \E(Y,A_i)\simeq k,\qquad \E(P_j,A_i)\simeq k\qquad(j<i).
\end{equation}
The category $\C(D)$ has enough injectives, and $\E^n=0$ for all $n\ge 2$.
\end{proposition}
\begin{proof}
We first identify the indecomposables and then compute their extensions using
the interval resolutions.

\medskip\noindent\textbf{Step 1: Indecomposables and idempotents.}
By \eqref{eq:hom} and \eqref{eq:X}, each listed object has endomorphism ring
$k$, and for any two distinct objects on the list at least one of the two
morphism spaces between them vanishes; hence they are pairwise nonisomorphic
indecomposables. Since every object of $\C(D)$ is a finite direct sum of these
objects and their endomorphism rings are local, $\C(D)$ is Krull--Schmidt and,
by Krause \cite[Corollary~4.4]{Kr}, idempotent complete.

\medskip\noindent\textbf{Step 2: Extensions, injectives, and heredity.}
Since $\rho X=0$, we have $\E(X,-)=\E(-,X)=0$. By Lemma~\ref{lem:intervals},
$Y$ and the $P_i$ are injective in $\B_I$, so $\E(-,Y)=\E(-,P_i)=0$ by
Construction~\ref{con:category}. For $B\in\B_I$, applying $\B_I(B,-)$ to the
injective resolution $0\to A_i\xrightarrow{b_i}P_i\xrightarrow{c_i}Y\to0$ of
Lemma~\ref{lem:intervals} yields
\[
  \E(B,A_i)=\Ext^1_{\B_I}(B,A_i)\simeq\coker\bigl(\B_I(B,P_i)
                  \xrightarrow{(c_i)_*}\B_I(B,Y)\bigr).
\]
By \eqref{eq:hom}, if $B=A_j$ then $\B_I(A_j,Y)=0$; if $B=Y$ then
$\B_I(Y,P_i)=0$ and $\B_I(Y,Y)=k$; and if $B=P_j$ then $\B_I(P_j,Y)=kc_j$,
whereas $(c_i)_*(\B_I(P_j,P_i))$ equals $kc_i p_{ji}=kc_j$ for $i\le j$ and $0$
for $j<i$. This proves \eqref{eq:extensions}.

In particular, $X$, $Y$, and all $P_i$ are injective in $\C(D)$. Each $A_i$
occurs as the first term of the conflation $A_i\xrightarrow{b_i}P_i\xrightarrow{c_i}Y$
with injective middle term, and each injective object $U$ occurs as the first
term of the split conflation $U\xrightarrow{1_U}U\to0$. Taking finite direct
sums shows that $\C(D)$ has enough injectives.

Finally, by Gorsky, Nakaoka and Palu \cite[Definition~3.1 and Remark~3.2]{GNP},
the group $\E^2(C,A)$ is generated by products $\alpha\otimes\beta$ with
$\alpha\in\E(M,A)$ and $\beta\in\E(C,M)$ for $M\in\C(D)$. Writing
$M=\bigoplus_{r=1}^t M_r$ as a finite direct sum of indecomposables with
inclusions $\iota_r:M_r\to M$ and projections $\pi_r:M\to M_r$, we have
$1_M=\sum_{r=1}^t\iota_r\pi_r$, and therefore
\[
  \alpha\otimes\beta
  =\sum_{r=1}^t (\iota_r)^*\alpha\otimes(\pi_r)_*\beta
\]
with $(\iota_r)^*\alpha\in\E(M_r,A)$ and $(\pi_r)_*\beta\in\E(C,M_r)$. By
\eqref{eq:extensions}, $\E(M_r,A)$ can be nonzero only when $M_r$ is $Y$ or
some $P_j$, whereas $\E(C,M_r)$ can be nonzero only when $M_r$ is some $A_i$.
Since these two sets of indecomposables are disjoint, every term vanishes.
Hence $\E^2=0$, and all higher tensor powers $\E^n$ ($n\ge 2$) vanish as well.
\end{proof}

We now apply Theorem~\ref{thm:criterion} and Proposition~\ref{prop:size} to the
inverse system $D$ of Construction~\ref{con:tree}.

\begin{theorem}\label{thm:example}
Let $k$ be any field and let $D$ be the inverse system of
Construction~\ref{con:tree}. The essentially small $k$-linear extriangulated
category $\C(D)$ has the following properties.
\begin{enumerate}
\item The category $\C(D)$ has no triangulated embedding.
\item Every countable family of objects of $\C(D)$ lies in a full
extension-closed subcategory with an algebraic realization.
\item The category $\C(D)$ is hereditary and Krull--Schmidt, and has enough
injectives. It has $\aleph_1$ isomorphism classes of indecomposables, each with
endomorphism ring $k$. For all indecomposables $A,C$ we have
$\dim_k\E(C,A)\le1$, and every morphism space has at most countable dimension
over $k$.
\end{enumerate}
\end{theorem}
\begin{proof}
(1) By Construction~\ref{con:tree}, we have $\varprojlim D=0$ and $D_0=k$, so
the projection $\varprojlim D\to D_0$ is not surjective. Hence $\C(D)$ has no
triangulated embedding by $(1)\Rightarrow(2)$ of Theorem~\ref{thm:criterion}.

(2) Every object of $\C(D)$ is a finite direct sum of $X$, $Y$, and the
$A_i,P_i$ ($i<\omega_1$), so only countably many indices $i<\omega_1$ occur in
a given countable family of objects. Since every countable subset of $\omega_1$
is bounded in $\omega_1$, these indices all lie in $J=[0,m]$ for some
$m<\omega_1$. Then $\C_J$ contains the given family, is extension-closed in
$\C(D)$, and admits an algebraic realization by Corollary~\ref{cor:axioms}
because $m=\max J$.

(3) By Proposition~\ref{prop:size}, $\C(D)$ is Krull--Schmidt, has enough
injectives, and is hereditary (since $\E^2=0$). Each indecomposable has
endomorphism ring $k$, and $\dim_k\E(C,A)\le1$ for all indecomposables $A,C$ by
\eqref{eq:extensions}. The indecomposables consist of the two families
$\{A_i\}_{i<\omega_1}$ and $\{P_i\}_{i<\omega_1}$ together with $X$ and $Y$, so
there are $\aleph_1$ isomorphism classes of indecomposables. By \eqref{eq:hom}
and \eqref{eq:X}, the only morphism spaces between indecomposables that can
have dimension greater than $1$ are the spaces $\C(D)(X,A_i)=D_i=k^{(N_i)}$,
each of which has dimension $|N_i|\le\aleph_0$ by Construction~\ref{con:tree}.
Taking finite direct sums preserves this bound, so every morphism space in
$\C(D)$ has at most countable dimension over $k$.
\end{proof}

In contrast to having enough injectives, the category $\C(D)$ of
Theorem~\ref{thm:example} does not have enough projectives: in fact, there is
no conflation ending in $Y$ with projective middle term. Indeed, an object $U$
is projective precisely when $\E(U,-)=0$ \cite[Proposition~3.24]{NP}. Since
$\omega_1$ has no largest element, for each $j<\omega_1$ there exists $i>j$
with $\E(P_j,A_i)\simeq k$, and $\E(Y,A_0)\simeq k$; hence by
\eqref{eq:extensions}, the only indecomposable projective objects are $X$ and
the $A_i$. Every morphism from $X$ or $A_i$ to $Y$ is zero, whereas by
Construction~\ref{con:category} the functor $\rho$ maps the second morphism of
any conflation ending in $Y$ to an epimorphism onto $Y$. On the other hand,
because $\C(D)$ has enough injectives, every object admits an inflation into an
injective object, so $\C(D)$ has enough injective morphisms; by Gorsky, Nakaoka
and Palu \cite[Corollary~5.22]{GNP}, the contravariant left derived functors of
$\Hom$ therefore equip $\C(D)$ with one-sided negative extensions
$\E_{\mathrm{II}}^{-n}$. Thus $\C(D)$ admits these one-sided negative
extensions even though it has no triangulated embedding; the obstruction in
Theorem~\ref{thm:criterion} arises from the degree $-1$ morphisms supplied by a
triangulated target.

Can either the index set or the morphism spaces in Theorem~\ref{thm:example} be
made smaller? The following proposition shows that neither bound can be
improved within the family of categories $\C(D)$ of
Construction~\ref{con:category}.

\begin{proposition}\label{prop:bounds}
Let $D=(D_i)_{i\in I}$ be an inverse system as in
Construction~\ref{con:category}. The category $\C(D)$ has an algebraic
realization under either of the following hypotheses.
\begin{enumerate}
\item The ordinal $I$ has a countable cofinal subset.
\item Every $D_i$ is finite-dimensional.
\end{enumerate}
Consequently, if $\C(D)$ has no triangulated embedding, then $I$ has no
countable cofinal subset and some $D_i$ is infinite-dimensional; in particular,
$|I|\ge\aleph_1$ and $\dim_k\C(D)(X,A_i)\ge\aleph_0$ for some $i\in I$. For the
system $D$ of Construction~\ref{con:tree}, we have
\[
  |I|=\aleph_1,\qquad \max_{i\in I}\dim_k D_i=\aleph_0 .
\]
\end{proposition}
\begin{proof}
By Theorem~\ref{thm:criterion}, to show that $\C(D)$ has an algebraic
realization under (1) or (2), it suffices to extend any given $i\in I$ and
$u\in D_i$ to a compatible family in $\varprojlim D$.

(1) If $I$ has a largest element $m$, we choose $w\in D_m$ with $\pi_{mi}(w)=u$
and take $(\pi_{mj}(w))_{j\in I}$. Otherwise, we apply the surjective case of
the Mittag--Leffler argument via the following recursive lift due to Tukey, as
recorded by Henkin \cite[p.~224]{Hen}. Since $I$ has a countable cofinal subset
and no largest element, there exists a strictly increasing cofinal sequence
$i=i_0<i_1<\dots$ in $I$. Setting $u_{i_0}=u$, we use surjectivity of the
transition maps to choose recursively $u_{i_{n+1}}\in D_{i_{n+1}}$ with
$\pi_{i_{n+1},i_n}(u_{i_{n+1}})=u_{i_n}$ for all $n\ge 0$. For an arbitrary
$j\in I$, choose $n$ with $i_n\ge j$ and define $u_j=\pi_{i_n,j}(u_{i_n})$. If
$i_m\ge i_n\ge j$, then
$\pi_{i_m,j}(u_{i_m})=\pi_{i_n,j}\pi_{i_m,i_n}(u_{i_m})=\pi_{i_n,j}(u_{i_n})$,
so $u_j$ is independent of $n$; choosing $i_n\ge k$ for $j\le k$ similarly
gives $\pi_{kj}(u_k)=u_j$. Thus $(u_j)_{j\in I}\in\varprojlim D$ extends $u$.

(2) Surjectivity of the projections for finite-dimensional vector spaces over
an arbitrary field follows from compactness in the coset topology (see Ghandour
\cite[Example~2.19 and Theorem~2.21(2)]{Gh}); we give a direct proof by
duality. Because the transition maps $\pi_{ji}:D_j\to D_i$ are surjective, the
dual maps $\pi_{ji}^*:D_i^*\to D_j^*$ between the dual spaces
$D_i^*=\Hom_k(D_i,k)$ are injective, so every $D_i^*$ embeds into the direct
limit $Z=\varinjlim_{j\in I} D_j^*$. Given $u\in D_i$, we extend evaluation at
$u$ from $D_i^*$ to a linear functional $\varphi\in Z^*$. For each $j\in I$,
since $D_j$ is finite-dimensional, the restriction $\varphi|_{D_j^*}$ is
evaluation at a unique vector $u_j\in D_j$, with $u_i=u$. For $h\le j$ in $I$
and every $\alpha\in D_h^*$, we have
\[
  \alpha(\pi_{jh}(u_j))=(\pi_{jh}^*\alpha)(u_j)
  =\varphi(\pi_{jh}^*\alpha)=\varphi(\alpha)=\alpha(u_h),
\]
whence $\pi_{jh}(u_j)=u_h$. Hence $(u_j)_{j\in I}\in\varprojlim D$ extends $u$.

Now suppose that $\C(D)$ has no triangulated embedding. Then it has no
algebraic realization, so by (1) the ordinal $I$ has no countable cofinal
subset (in particular, $|I|\ge\aleph_1$), and by (2) some
$D_i=\C(D)(X,A_i)$ is infinite-dimensional (so $\dim_k\C(D)(X,A_i)\ge\aleph_0$).
For the system $D$ of Construction~\ref{con:tree}, we have $I=\omega_1$ (hence
$|I|=\aleph_1$) and $\dim_k D_i\le\aleph_0$ for all $i<\omega_1$; since
$\C(D)$ has no triangulated embedding by Theorem~\ref{thm:example}(1), some
$D_i$ is infinite-dimensional, and therefore $\max_{i\in I}\dim_k D_i=\aleph_0$.
\end{proof}

\begin{thebibliography}{Stacks}
\bibitem[AET]{AET}
T. Adachi, H. Enomoto and M. Tsukamoto, \emph{Intervals of $s$-torsion
pairs in extriangulated categories with negative first extensions},
Math. Proc. Cambridge Philos. Soc. \textbf{174} (2023), no.~3, 451--469.
\bibitem[Be]{Be}
G.~M. Bergman, \emph{Every module is an inverse limit of injectives},
Proc. Amer. Math. Soc. \textbf{141} (2013), no.~4, 1177--1183.
\bibitem[Be2]{Be2}
G.~M. Bergman, \emph{Some empty inverse limits}, July 6, 2011,
\url{https://math.berkeley.edu/~gbergman/papers/unpub/emptylim.pdf}.
\bibitem[BS]{BS}
R. Bennett-Tennenhaus and A. Shah, \emph{Transport of structure in higher
homological algebra}, J. Algebra \textbf{574} (2021), 514--549.
\bibitem[BCB]{BCB}
M.~B. Botnan and W. Crawley-Boevey, \emph{Decomposition of persistence
modules}, Proc. Amer. Math. Soc. \textbf{148} (2020), 4581--4596.
\bibitem[Ch]{Ch}
X. Chen, \emph{Exact dg categories II: The embedding theorem},
J. Pure Appl. Algebra \textbf{230} (2026), no.~7, article 108298.
\bibitem[Di]{Di}
R.~M. Dimitri\'c, \emph{Trees, inverse systems, and valuated vector
spaces}, Acta Math. Univ. Ostrav. \textbf{12} (2004), no.~1, 19--22.
\bibitem[Gh]{Gh}
N.~J. Ghandour, \emph{A note on the inverse limits of linear algebraic
groups}, Int. J. Algebra \textbf{5} (2011), no.~19, 925--933.
\bibitem[GNP]{GNP}
M. Gorsky, H. Nakaoka and Y. Palu, \emph{Positive and negative extensions
in extriangulated categories}, Math. Z. \textbf{313} (2026), article 89.
\bibitem[GNP2]{GNP2}
M. Gorsky, H. Nakaoka and Y. Palu, \emph{Hereditary extriangulated
categories: Silting objects, mutation, negative extensions},
\href{https://arxiv.org/abs/2303.07134v2}{arXiv:2303.07134v2} (2023).
\bibitem[Hen]{Hen}
L. Henkin, \emph{A problem on inverse mapping systems},
Proc. Amer. Math. Soc. \textbf{1} (1950), 224--225.
\bibitem[HeS]{HeS}
E. Herscovich and A. Solotar, \emph{Derived invariance of
Hochschild--Mitchell (co)homology and one-point extensions},
J. Algebra \textbf{315} (2007), no.~2, 852--873.
\bibitem[HS]{HS}
G. Higman and A.~H. Stone, \emph{On inverse systems with trivial limits},
J. London Math. Soc. (1) \textbf{29} (1954), no.~2, 233--236.
\bibitem[Je]{Je}
T. Jech, \emph{Set Theory}, third millennium edition, revised and expanded,
Springer Monographs in Mathematics, Springer-Verlag, Berlin, 2003.
\bibitem[Ke]{Ke}
B. Keller, \emph{Calabi--Yau triangulated categories},
in \emph{Trends in representation theory of algebras and related topics},
EMS Ser. Congr. Rep., Eur. Math. Soc., Z\"urich, 2008, 467--489.
\bibitem[Kl]{Kl}
J. Klemenc, \emph{The stable hull of an exact $\infty$-category},
Homology Homotopy Appl. \textbf{24} (2022), no.~2, 195--220.
\bibitem[Kr]{Kr}
H. Krause, \emph{Krull--Schmidt categories and projective covers},
Expo. Math. \textbf{33} (2015), no.~4, 535--549.
\bibitem[Ku]{Ku}
G. Kurepa, \emph{Ensembles ordonn\'es et ramifi\'es},
doctoral thesis, Universit\'e de Paris, 1935.
\bibitem[NP]{NP}
H. Nakaoka and Y. Palu, \emph{Extriangulated categories, Hovey twin cotorsion
pairs and model structures}, Cah. Topol. G\'eom. Diff\'er. Cat\'eg.
\textbf{60} (2019), no.~2, 117--193.
\bibitem[NP2]{NP2}
H. Nakaoka and Y. Palu, \emph{External triangulation of the homotopy
category of exact quasi-category},
\href{https://arxiv.org/abs/2004.02479v1}{arXiv:2004.02479v1} (2020).
\bibitem[Stacks]{Stacks}
The Stacks Project Authors, \emph{The Stacks Project},
\href{https://stacks.math.columbia.edu/tag/05TG}{Tag 05TG}.
\bibitem[Wa]{Wa}
W.~C. Waterhouse, \emph{An empty inverse limit},
Proc. Amer. Math. Soc. \textbf{36} (1972), 618.
\bibitem[Wo]{Wo}
E. Wofsey, \href{https://math.stackexchange.com/a/3303841}{Mathematics
Stack Exchange, answer 3303841}, July 25, 2019, revised July 27, 2019.
\end{thebibliography}
\end{document}
